%!TeX root=MemoriaTFG.tex

\chapter{Metodología}
\label{cap:metodos}

Este capítulo documenta la metodología del trabajo con el detalle
suficiente para que cualquier cifra del
capítulo~\ref{cap:resultados} pueda reconstruirse a partir de los
artefactos publicados. Se describen los datos empleados, los modelos
comparados y la infraestructura de cómputo, junto con las métricas,
las pruebas estadísticas y los experimentos que sostienen la
evaluación.

\section{Datasets}
\label{sec:datos}

La selección de datasets persigue medir la capacidad de
generalización de los modelos entre modalidades de imagen, dominios
clínicos y tareas. Por ello se emplean doce datasets externos
públicos que cubren dermatoscopia, fotografía clínica,
histopatología, análisis de equidad y aprendizaje basado en
conceptos, agrupados por modalidad y finalidad. La
tabla~\ref{tab:datasets-resumen} resume sus características.

\begin{table}[htbp]
\centering
\footnotesize
\setlength{\tabcolsep}{4pt}
\renewcommand{\arraystretch}{0.95}
\caption{Datasets externos empleados en la evaluación.
$N_{\mathrm{test}}$: tamaño nominal del \emph{split} de evaluación
canónico; ``Clases'': número de categorías diagnósticas o conceptos
en la convención propia del dataset. El asterisco
(\textsuperscript{$\ast$}) marca Dermnet, el único dataset con
solapamiento exacto frente al dataset de preentrenamiento Derm1M
(auditoría en sección~\ref{sec:auditoria-resultados}). Las dos variantes de Derm7pt
(clínica y dermatoscópica) comparten origen
(Seven-Point Checklist) y se reportan como un único dataset en el
conteo agregado de doce datasets externos.}
\label{tab:datasets-resumen}
\begin{tabular}{lrcll}
\hline
Dataset & $N_{\mathrm{test}}$ & Clases & Modalidad & Uso \\
\hline
HAM10000~\cite{ham10000}             & 1\,232  & 7        & Dermoscopia      & LP, FT, ZS, LLMs \\
BCN20000~\cite{bcn20000}             & 1\,242  & 9        & Dermoscopia      & LP \\
PAD-UFES-20~\cite{padufes}           & 461     & 6        & Clínica móvil    & LP, FT \\
DDI~\cite{ddi}                       & 137     & 2        & Clínica          & LP, FT, equidad \\
Dermnet~\cite{dermnet}               & 4\,002  & 23       & Clínica          & LP\textsuperscript{$\ast$} \\
Derm7pt clínico~\cite{kawahara2019}  & 168     & 2        & Clínica          & LP \\
Derm7pt dermo~\cite{kawahara2019}    & 225     & 2        & Dermoscopia      & LP \\
HIBA                                 & 334     & 2        & Dermoscopia      & LP \\
MSKCC                                & 1\,664  & 2        & Dermoscopia      & LP, FT \\
WSI patches                          & 12\,354 & 16       & Histopatología   & LP \\
ISIC2018~\cite{isic2018}             & 2\,594  & binaria  & Dermoscopia      & Segmentación \\
Fitzpatrick17k~\cite{fitzpatrick17k} & 16\,577 & 114 + FP & Clínica          & Equidad \\
SkinCon~\cite{skincon}               & 3\,230  & 48 conc. & Clínica          & CBM \\
\hline
\end{tabular}
\end{table}

Los datasets cubren las tres modalidades del diagnóstico
dermatológico. La dermatoscopia (HAM10000, BCN20000, Derm7pt dermo,
HIBA, MSKCC, ISIC2018) aporta imagen polarizada adquirida con
dermatoscopio. La fotografía clínica (PAD-UFES-20, DDI, Dermnet,
Derm7pt clínico) reproduce las condiciones de captura en consulta y
atención primaria. La histopatología la aporta WSI~patches, con
$12\,354$ \emph{parches} en $16$ categorías. Dos datasets cumplen
un papel específico: Fitzpatrick17k habilita el análisis de equidad,
al anotar conjuntamente patología ($114$ categorías) y fototipo
cutáneo (escala Fitzpatrick I-VI), y SkinCon aporta anotación densa
de $48$ conceptos clínicos visuales para los Concept Bottleneck
Models. En todos los experimentos se utilizaron las particiones oficiales de entrenamiento, validación y prueba \emph{train/val/test} proporcionadas por cada dataset, sin modificar ni reasignar imágenes entre ellas.


\subsection{Independencia frente al preentrenamiento: solapamiento con Derm1M}
\label{sec:auditoria-resultados}

La validez de toda comparación de transferencia y \emph{zero-shot}
del trabajo depende de que los conjuntos de test no se solapen con el
dataset de preentrenamiento, por lo que conviene comprobarlo antes de
emplearlos. La auditoría ---cuyo protocolo se detalla en la
sección~\ref{sec:auditoria}--- intersecta los hashes MD5 de los doce
datasets externos con los $415\,451$ hashes únicos de la fracción
accesible de Derm1M e identifica un único dataset con solapamiento
exacto: \textbf{Dermnet}, cuyas $18\,302$ imágenes están todas
contenidas en Derm1M ($100\%$). Los once datasets restantes no
arrojan ninguna coincidencia. En particular, HIBA y MSKCC ---cuyo
origen en el ISIC archive motivaba una sospecha de procedencia---
presentan $0$ coincidencias MD5, por lo que esa sospecha no se
traduce en solapamiento exacto. El resultado completo por dataset
figura en el repositorio reproducible (\ref{cap:repo},
sección~\ref{sec:repo-reproducibilidad}).

El solapamiento de Dermnet no invalida las cifras que se reportan en
la sección~\ref{sec:rep-panderm} para PanDerm, DINOv2,
ConvNeXt-Large, EfficientNetV2-Large y SigLIP-Large, cuyo
preentrenamiento no incluye Derm1M. Sí condiciona, en cambio, la
interpretación de las cifras zero-shot de la familia DermLIP sobre
Dermnet (sección~\ref{sec:zs-results}), ya que DermLIP se entrenó
sobre Derm1M y vio durante el preentrenamiento la totalidad de ese
conjunto.

\section{Ontología jerárquica L1/L2/L3}
\label{sec:ontologia}

Comparar modelos entre datasets exige resolver un problema previo:
cada conjunto usa su propia taxonomía, y esas taxonomías no son
compatibles entre sí. Dos diagnósticos equivalentes pueden aparecer
bajo denominaciones distintas o con niveles de granularidad
diferentes. Para unificarlas se diseñó, en colaboración
con la Dra.~R.~Taberner, una ontología jerárquica de tres niveles
que actúa como vocabulario común a todos los protocolos del trabajo.

El nivel L1 (etiológico) reparte los diagnósticos en cuatro
familias: patología inflamatoria, tumoral, infecciosa y
genodermatosis. El nivel L2 (subcategoría clínica) contiene
$43$ entradas; cinco de ellas no se desagregan en L3 por
considerarse ya suficientemente específicas (\emph{Esclerosis
tuberosa}, \emph{Neurofibromatosis tipo~1}, \emph{Incontinentia
pigmenti}, \emph{Queratodermia palmoplantar hereditaria} y
\emph{Pseudoxantoma elástico}). El nivel L3 (diagnóstico) contiene
$367$ entradas individuales, cada una con sus códigos SNOMED~CT y
CIE-10 para garantizar interoperabilidad con sistemas clínicos.

Once de los doce datasets se mapean a la ontología; SkinCon queda
fuera por estar anotado con conceptos, no con diagnósticos. El
conjunto armonizado resultante reúne $72\,654$ imágenes etiquetadas
sobre vocabulario común.

\section{Dataset DermapixelAI 1.0}
\label{sec:dermapixel}

DermapixelAI~1.0 es un dataset dermatológico en español construido
en este trabajo a partir del archivo del blog Dermapixel de la
Dra.~R.~Taberner; su construcción materializa el objetivo~O2
(sección~\ref{sec:objetivos}). A diferencia de los datasets públicos
utilizados en la evaluación, DermapixelAI~1.0 incorpora texto
narrativo clínico en español asociado a cada caso, lo que permite
estudiar escenarios multimodales y de transferencia lingüística no
cubiertos por los conjuntos de referencia internacionales. La
elaboración comprende la
extracción del material del archivo, la depuración, la
estructuración que vincula para cada caso imagen y modalidad,
metadatos clínicos y texto narrativo, y el etiquetado contra la
ontología jerárquica L1/L2/L3; el pipeline completo se documenta en
el repositorio reproducible (\ref{cap:repo}, sección~\ref{sec:repo-ablaciones}). La versión utilizada reúne $1\,089$
imágenes vinculadas a $669$ casos clínicos, cada uno con su texto
narrativo asociado (mediana de $223$ palabras). Predomina la
fotografía clínica ($1\,036$ imágenes), con un subconjunto
dermatoscópico reducido ($49$).

Las imágenes se etiquetan según la ontología L1/L2/L3 y se
particionan bajo regla \emph{case-aware} ($891 / 157 / 41$
imágenes; $538 / 100 / 31$ casos), de modo que todas las imágenes
de un mismo caso caen en el mismo \emph{split} y no hay fuga entre
conjuntos. Dos campos cuantifican la calidad del etiquetado: el
$97{,}89\%$ del dataset está mapeado a la ontología con revisión
experta del vocabulario (\texttt{label\_source = ontology}) y el
$84{,}39\%$ ha pasado además revisión visual caso-a-caso por la
Dra.~R.~Taberner (\texttt{rosa\_verified}). Ambas cifras no son
intercambiables y se detallan, junto a un tercer campo de validación
textual, en el repositorio reproducible (\ref{cap:repo}, sección~\ref{sec:repo-ablaciones}).

El dataset se distribuye bajo licencia CC-BY-NC-SA~4.0. Su pipeline
de construcción, caracterización completa y datasheet figuran en el
repositorio reproducible (\ref{cap:repo}, sección~\ref{sec:repo-ablaciones}), y el documento formal de entrega, en su
sección~\ref{sec:repo-datasheet}. La validación experimental sobre
DermapixelAI~1.0 forma parte de los resultados de este trabajo: las
secciones~\ref{sec:dermapixel-results} y~\ref{sec:dermapixel-zs} del
capítulo~\ref{cap:resultados} evalúan sobre este dataset el sondeo
lineal supervisado y la clasificación zero-shot multimodal en los
tres niveles ontológicos.

\section{Modelos evaluados}
\label{sec:modelos}

Los modelos incluidos en el estudio se seleccionan con el objetivo
de representar los principales paradigmas actualmente utilizados en
visión artificial y aprendizaje multimodal aplicado a dermatología.
La comparación combina modelos específicamente entrenados para
dermatología, modelos visuales generalistas de referencia,
arquitecturas visión-lenguaje biomédicas y modelos multimodales de
propósito general. Esta diversidad permite evaluar hasta qué punto
la especialización dermatológica aporta ventajas medibles frente a
modelos más generales y situar el rendimiento de los modelos
fundacionales dermatológicos dentro del ecosistema actual de
inteligencia artificial médica.

La tabla~\ref{tab:modelos} resume los dieciséis sistemas evaluados:
catorce modelos fundacionales ---encoders visuales de dominio
(PanDerm) y generales (DINOv2), modelos visión-lenguaje (la familia
DermLIP, BiomedCLIP y SigLIP), el segmentador promptable SAM2.1 y
los modelos de lenguaje multimodales (GPT-4o, Gemini~2.5~Pro,
MedGemma y BLIP-2)--- y dos líneas de base de redes convolucionales
supervisadas (ConvNeXt-Large y EfficientNetV2-Large), incluidas como
referencia no fundacional frente a la que medir el valor del
preentrenamiento.

\begin{table}[htbp]
\centering
\footnotesize
\setlength{\tabcolsep}{4pt}
\renewcommand{\arraystretch}{0.95}
\caption{Modelos evaluados.}
\label{tab:modelos}
\begin{tabular}{llrl}
\hline
Modelo & Familia & Parámetros & Referencia \\
\hline
PanDerm Base                & Encoder visual derm.       & $\sim 86$\,M & \cite{panderm2025} \\
PanDerm Large               & Encoder visual derm.       & $\sim 307$\,M & \cite{panderm2025} \\
DINOv2 ViT-L/14             & Encoder visual general     & $\sim 300$\,M & \cite{dinov2} \\
ConvNeXt-Large              & ConvNet supervisada        & $\sim 198$\,M & --- \\
EfficientNetV2-Large        & ConvNet supervisada        & $\sim 118$\,M & --- \\
DermLIP v1 (PubMedBERT)     & Vision-lenguaje derm.      & --- & \cite{derm1m} \\
DermLIP v2 (PubMedBERT)     & Vision-lenguaje derm.      & --- & \cite{derm1m} \\
DermLIP original (GPT-2)    & Vision-lenguaje derm.      & --- & \cite{derm1m} \\
BiomedCLIP                  & Vision-lenguaje biomédico  & --- & \cite{biomedclip} \\
SigLIP-Large SO400M         & Vision-lenguaje general    & $\sim 878$\,M & \cite{siglip2023} \\
SAM2.1-Large (decoder FT)   & Segmentación               & $\sim 227$\,M & \cite{sam2} \\
GPT-4o                      & LLM multimodal             & no público & --- \\
Gemini 2.5 Pro              & LLM multimodal             & no público & --- \\
MedGemma 4B Vision          & LLM multimodal biomédico   & $\sim 4$\,B & --- \\
MedGemma 27B Text+LoRA      & LLM multimodal biomédico   & $\sim 27$\,B & --- \\
BLIP-2 (Q-Former)           & Vision-lenguaje general    & --- & \cite{blip2_2023} \\
\hline
\end{tabular}
\end{table}

Los encoders difieren en la dimensión de sus \emph{embeddings}:
$768$ en PanDerm Base; $1\,024$ en PanDerm Large y DINOv2
ViT-L/14; $1\,152$ en SigLIP-Large SO400M. La cifra de
$\sim 307$\,M parámetros de PanDerm Large es el valor nominal del
ViT-L/16 publicado por sus autores. Al cargarlo como extractor de
características y sustituir la cabeza original por la identidad
quedan $303{,}3$\,M parámetros efectivos. Esta diferencia no afecta
al rendimiento, porque la cabeza descartada no interviene en la
inferencia. La familia DermLIP usa ViT-B/16 como encoder visual y
dos encoders textuales alternativos: PubMedBERT extendido a $256$
tokens (variantes v1 y v2) y GPT-2 con tokenizador de CLIP a $77$
tokens (variante original). DermFM-Zero~\cite{dermfmzero2026} se
incluye en la discusión pero no en la evaluación empírica, ya que
sus pesos no estaban disponibles a la fecha de cierre.

\section{Infraestructura experimental}
\label{sec:infraestructura}

Todos los experimentos corren sobre una NVIDIA DGX Spark con chip
Grace Hopper GB10, $128$\,GB de memoria unificada compartida entre
CPU y GPU, Linux ARM aarch64 y CUDA~13.0. La precisión nativa es
BF16 (\emph{brain floating-point} de 16 bits) en inferencia y
entrenamiento, con FP32 en optimizador y pérdida durante el ajuste
fino (precisión mixta canónica). La memoria unificada es clave para
este trabajo: permite mantener cargados modelos del orden de
$55$\,GB en BF16 sin \emph{offloading} ni cuantización, de modo
que incluso MedGemma~27B opera localmente y sin comprimir. El
software queda fijado en Python~3.12.3, PyTorch~2.11.0,
\texttt{timm}~0.9.16, \texttt{scikit-learn}, \texttt{transformers},
\texttt{open\_clip}, \texttt{sam2}, \texttt{peft} y
FAISS~\cite{faiss2019}.

Los tiempos de cálculo característicos sobre este equipo orientan
sobre el coste de cada protocolo: un sondeo lineal completo de un
encoder sobre un dataset estándar tarda $1$--$3$ minutos; el
ajuste fino de PanDerm Large sobre HAM10000 ($50$ épocas con TTA),
unas $15$ horas; el ajuste fino del decodificador de SAM2.1-Large
sobre ISIC2018 ($50$ épocas), unas $24$ horas; la validación
cruzada de cinco pliegues del modelo de segmentación, con la
ablación de rango \emph{LoRA} asociada, unas $18$ horas; y el
entrenamiento de un Sparse Autoencoder Large ($16\,384$
características), unas $4$ horas.

\section{Métricas}
\label{sec:metricas}

Ninguna métrica individual captura por sí sola el rendimiento sobre
problemas dermatológicos, caracterizados por clases desbalanceadas,
distintos niveles de granularidad diagnóstica y requisitos clínicos
diversos. El trabajo reporta, por tanto, un conjunto de métricas
complementarias para clasificación, segmentación, equidad y
rendimiento computacional.

\paragraph{Clasificación.}
La métrica principal es la \emph{Balanced Accuracy} (BAcc, promedio
del \emph{recall} por clase): se adopta porque asigna el mismo peso
a todas las clases con independencia de su frecuencia, lo que evita
que las categorías mayoritarias dominen la evaluación en los
datasets desbalanceados de este trabajo. La \emph{accuracy} acompaña
como referencia comparativa
entre modelos sobre la misma distribución. Se reportan además la F1
ponderada (media armónica de precisión y \emph{recall}, pesada por
frecuencia de clase), la AUROC \emph{one-vs-rest}
(\texttt{multi\_class='ovr'} en \texttt{scikit-learn}, promediada
sobre clases) y el coeficiente kappa de Cohen, que corrige la
concordancia por azar. \emph{Recall@$k$} —presencia de la etiqueta
verdadera entre las $k$ primeras posiciones del ranking— se reporta
específicamente para melanoma sobre HAM10000.

\paragraph{Segmentación.}
La calidad de las máscaras predichas se evalúa mediante el
coeficiente Dice (DSC) y la intersección sobre unión (IoU), dos
métricas estándar que cuantifican el grado de solapamiento entre la
segmentación automática y la máscara de referencia:
$\mathrm{DSC} = 2\,|P\cap G|/(|P|+|G|)$ e
$\mathrm{IoU} = |P\cap G|/|P\cup G|$, con $P$ la máscara predicha y
$G$ la de referencia.

\paragraph{Equidad.}
BAcc por subgrupo Fitzpatrick (I-VI) sobre la formulación de tres
clases de malignidad, y \emph{gap}
$\mathrm{Gap}_{\mathrm{I}\to\mathrm{VI}}
= \mathrm{BAcc}_{\mathrm{I}} - \mathrm{BAcc}_{\mathrm{VI}}$. Un
valor próximo a cero indica rendimiento comparable entre fototipos
extremos; su interpretación conjunta con la BAcc global se discute en
la sección~\ref{sec:equidad}.

\paragraph{Operativas.}
Además del rendimiento predictivo, se evalúa el coste computacional
asociado a cada modelo, que permite valorar la viabilidad de un
despliegue en entornos clínicos con recursos limitados: latencia por
consulta en milisegundos (inferencia más posproceso) y huella de
memoria (VRAM en gigabytes, BF16), medidas sobre la infraestructura
de sección~\ref{sec:infraestructura}.

\paragraph{Agregación entre datasets.}
\label{sec:agregacion-datasets}
Algunas vistas panorámicas del capítulo~\ref{cap:resultados} reportan
la media de una métrica sobre varios datasets externos. Esta media
debe interpretarse únicamente como un resumen descriptivo de
tendencias generales entre conjuntos con taxonomías, modalidades y
prevalencias diferentes. No pretende establecer un ranking absoluto
entre modelos, ya que datasets como HAM10000, DDI, PAD-UFES-20 o
Dermnet no son estrictamente comparables entre sí. Por ello, todas las
conclusiones se fundamentan en el análisis desglosado por dataset y
por nivel ontológico, utilizando las métricas agregadas únicamente
como apoyo descriptivo.

\section{Pruebas estadísticas}
\label{sec:estadistica}

Que un modelo obtenga una cifra mejor que otro no basta para afirmar
que sea mejor: la diferencia podría deberse al azar de qué imágenes
concretas cayeron en el conjunto de evaluación. Este riesgo es mayor
en dermatología, donde los conjuntos de test son moderados, las
clases están desbalanceadas y los casos son clínicamente
heterogéneos. Para no confundir una diferencia real con una
casualidad, el trabajo adopta un protocolo estadístico explícito,
alineado con las guías de reporte para inteligencia artificial en
imagen médica (CLAIM~\cite{claim2024} y TRIPOD+AI~\cite{tripodai2024}).
Lo componen tres herramientas: los intervalos de confianza por
\emph{bootstrap}, que miden cuánto puede oscilar cada cifra; el test
de DeLong, que compara dos modelos de forma justa; y la corrección de
Holm, que evita dar por buenas diferencias surgidas de hacer muchas
comparaciones.

El \emph{bootstrap} responde a la pregunta ``¿cuánto cambiaría esta
cifra si el test hubiera estado formado por otras imágenes?''. El
procedimiento es sencillo: se generan $1\,000$ conjuntos de test
``alternativos'' tomando imágenes del test real al azar y con
reemplazo ---conservando la proporción de clases---, y se recalcula la
métrica en cada uno. La dispersión de esas $1\,000$ cifras describe su
incertidumbre, y el rango que descarta el $2{,}5\%$ de valores más
altos y el $2{,}5\%$ más bajos define el intervalo de confianza al
$95\%$: el margen dentro del cual cabe esperar el valor verdadero.
Este método no exige suponer que los datos sigan ninguna distribución
concreta, lo que lo hace apropiado para los conjuntos reducidos y
desbalanceados de este trabajo, donde las fórmulas estadísticas
clásicas no son fiables. Los intervalos al $95\%$ ($B = 1\,000$,
remuestreo estratificado por clase, percentil $2{,}5/97{,}5$) de las
métricas \emph{headline} que no los incluyen en el cuerpo figuran en
el repositorio reproducible (\ref{cap:repo},
sección~\ref{sec:repo-bootstrap-ci}).

Para decidir si la diferencia de AUROC entre dos modelos es real, y no
fruto del azar, se emplea el test de DeLong~\cite{delong1988}. La
razón de elegirlo es importante: como ambos modelos se evalúan sobre
exactamente las mismas imágenes, sus aciertos y errores están ligados;
tratar sus dos AUROC como si procedieran de muestras independientes
haría parecer las diferencias más significativas de lo que son. El
test de DeLong tiene en cuenta esa relación y produce una comparación
pareada correcta. Se aplica a los contrastes entre codificadores sobre
tres tareas binarias de detección de malignidad ---Fitzpatrick17k,
HAM10000 y DermapixelAI~1.0--- (sección~\ref{sec:equidad-delong};
tabla detallada en el repositorio reproducible, \ref{sec:repo-tablas}).

Por último, cuando se hacen muchas comparaciones aumenta la
probabilidad de que alguna parezca significativa solo por azar. Para
controlarlo, los $p$-valores se reportan de dos maneras: el valor
nominal (sin corregir) y el valor corregido por el método de Holm
dentro de cada tarea. Mostrar ambos es deliberado: la corrección
reduce los falsos positivos, pero también la capacidad de detectar
diferencias reales, y ofrecer las dos cifras permite al lector juzgar
por sí mismo la solidez de cada comparación. En la práctica, las
diferencias observadas son tan grandes (con $p$ varios órdenes de
magnitud por debajo de $0{,}05$) que las conclusiones principales se
mantienen con cualquier corrección estándar. Los tamaños de muestra
por subgrupo ---por fototipo cutáneo y por nivel de la ontología--- se
recogen en las tablas~\ref{tab:muestral-fst} y~\ref{tab:muestral-nivel},
que ayudan a ver qué análisis tienen más respaldo: en particular, el
escaso número de casos del fototipo~VI ($n = 73$ en test) y la larga
cola de diagnósticos poco frecuentes de L3 limitan la fiabilidad de
las conclusiones en esos subgrupos. El detalle por subgrupo acompaña
además a sus cifras en el capítulo~\ref{cap:resultados}.

\begin{table}[htbp]
\centering
\footnotesize
\caption{Tamaño muestral por fototipo cutáneo (Fitzpatrick17k),
cohorte de los análisis de equidad y del test de DeLong de
malignidad. La evaluación del \emph{gap}~I-VI usa los $1\,597$ casos
de test con fototipo asignado; el \emph{endpoint} binario de
malignidad emplea los $1\,658$ del test ($226$ malignos). El reducido
tamaño del fototipo~VI ($n = 73$ en test) acota la potencia
estadística en ese subgrupo.}
\label{tab:muestral-fst}
\begin{tabular}{lrr}
\hline
Fototipo (FST) & $N$ corpus & $N$ test \\
\hline
I    & $2\,947$ & $299$ \\
II   & $4\,808$ & $470$ \\
III  & $3\,308$ & $325$ \\
IV   & $2\,781$ & $261$ \\
V    & $1\,533$ & $169$ \\
VI   & $635$    & $73$  \\
Sin asignar & $565$ & $61$ \\
\hline
\textbf{Total} & $\mathbf{16\,577}$ & $\mathbf{1\,658}$ \\
\hline
\end{tabular}
\end{table}

\begin{table}[htbp]
\centering
\footnotesize
\caption{Tamaño muestral por nivel ontológico sobre DermapixelAI~1.0
($1\,089$ imágenes, $669$ casos). El vocabulario común declara $4$
categorías etiológicas L1, $43$ subcategorías clínicas L2 y $367$
diagnósticos L3; el dataset armonizado de los doce externos aporta
$72\,654$ imágenes mapeadas a ese vocabulario. La cola larga L3
concentra $64$ clases con una sola imagen, $129$ con dos a cinco, $39$
con seis a diez, $11$ con once a veinte y $7$ con veintiuna o más.}
\label{tab:muestral-nivel}
\begin{tabular}{lrr}
\hline
Nivel & Declarado & Representado \\
\hline
L1 & $4$   & $4$ \\
L2 & $43$  & $38$ \\
L3 & $367$ & $250$ \\
\hline
\end{tabular}
\end{table}

Conviene aclarar que este aparato estadístico no se aplica por igual
en todo el trabajo. El test de DeLong y la corrección por
comparaciones múltiples se usan allí donde comparar modelos es el
objetivo central ---el ranking de codificadores en detección de
malignidad (sección~\ref{sec:equidad-delong})---, y los intervalos por
\emph{bootstrap} acompañan a las métricas principales sobre
DermapixelAI~1.0 y a la validación cruzada de segmentación. Otros
bloques ---el sondeo lineal sobre los datasets externos, la eficiencia
en etiquetas y la equidad por fototipo--- se presentan como una sola
estimación, sin repetirla con varias semillas
(sección~\ref{sec:lim-semilla}). Por eso las afirmaciones se gradúan
según el respaldo estadístico disponible: se presentan como patrones
sólidos cuando la diferencia supera con holgura la variabilidad
esperable, y con la cautela correspondiente cuando no se dispone de un
contraste formal.

\section{Trazabilidad y reproducibilidad}
\label{sec:reproducibilidad}

La reproducibilidad es una condición necesaria para la validación
independiente de los resultados. El trabajo se diseña, por tanto,
para que cualquier número del capítulo~\ref{cap:resultados} se
reconstruya a partir de artefactos persistentes y no de la memoria
de una ejecución concreta. Tres mecanismos articulan esta garantía.

El primero es el determinismo de la ejecución. Las semillas
aleatorias se fijan a $0$ para \texttt{torch}, \texttt{numpy} y
\texttt{random} en \emph{fine-tuning} y segmentación, y a $42$ en
\texttt{random\_state} para los protocolos basados en
\texttt{scikit-learn} (sondeo lineal, eficiencia en etiquetas,
equidad). El entorno queda fijado a las versiones de
sección~\ref{sec:infraestructura}, y cualquier desviación se
documenta en el experimento afectado. Por último, se deshabilita la
telemetría externa y se fija explícitamente el dispositivo de
ejecución para evitar variaciones no controladas del entorno
experimental.

El segundo mecanismo es la persistencia de las salidas. Los
resultados numéricos se vuelcan a ficheros CSV o JSON con esquema
estable, lo que permite reconstruir cualquier tabla o figura del
capítulo~\ref{cap:resultados} sin reejecutar nada, y los mapeos de
etiquetas nativas a la ontología L1/L2/L3 se mantienen fijos entre
experimentos. El tercero es la integridad de la cadena de cómputo:
cuando un experimento toma como entrada la salida de otro ---por
ejemplo, los \emph{embeddings} del sondeo lineal alimentan la rama
visual del zero-shot---, la identidad de esa entrada se preserva
mediante hash MD5 del fichero. El encadenamiento por hash convierte
el conjunto de experimentos en un grafo de dependencias verificable:
cualquier alteración silenciosa de una entrada intermedia
invalidaría el hash y sería detectable, lo que protege la cadena de
evidencia frente a la deriva de datos entre etapas dependientes. El
listado completo de tareas reproducibles, con sus comandos, ficheros
de salida y \emph{checkpoints}, figura en el
repositorio reproducible (\ref{cap:repo}, sección~\ref{sec:repo-reproducibilidad}).

\section{Experimentos}
\label{sec:protocolos}

Los experimentos se organizan en cuatro bloques. El primero mide la
calidad y la adaptabilidad de las representaciones visuales: el
sondeo lineal evalúa lo que capturan los encoders congelados, el
ajuste fino su rendimiento máximo tras adaptación supervisada, y la
eficiencia en etiquetas cuántos datos anotados hacen falta para
alcanzarlo. El segundo aborda la localización espacial de la lesión
mediante segmentación. El tercero examina la capacidad multimodal: la
clasificación \emph{zero-shot} guiada por texto y la comparación con
los grandes modelos de lenguaje generalistas. El cuarto protege la
validez de las conclusiones: la evaluación de equidad analiza el
rendimiento entre fototipos cutáneos y la auditoría de solapamiento
con Derm1M descarta la fuga de información entre evaluación y
preentrenamiento. La tabla~\ref{tab:experimentos} resume cada
experimento, la pregunta que responde y su alcance. Todos comparten
la infraestructura (sección~\ref{sec:infraestructura}) y las semillas
declaradas (sección~\ref{sec:reproducibilidad}), de modo que sus
cifras son comparables entre sí.

\begin{table}[htbp]
\centering
\footnotesize
\setlength{\tabcolsep}{4pt}
\renewcommand{\arraystretch}{1.15}
\caption{Resumen de los experimentos: la pregunta que responde cada
uno y su alcance en modelos y datasets. DermapixelAI~1.0 se evalúa
además en sondeo lineal y \emph{zero-shot}
(secciones~\ref{sec:dermapixel-results} y~\ref{sec:dermapixel-zs}).}
\label{tab:experimentos}
\begin{tabular}{p{2.5cm}p{3.7cm}p{2.9cm}p{3.4cm}}
\hline
Experimento & Pregunta & Modelos & Datasets \\
\hline
Sondeo lineal (\ref{sec:lp})
  & Calidad de las representaciones congeladas
  & $8$ encoders visuales
  & $10$ externos de clasificación \\
Ajuste fino (\ref{sec:ft})
  & Rendimiento máximo tras adaptación supervisada
  & PanDerm Base/Large
  & HAM10000, PAD-UFES-20, DDI, MSKCC \\
Eficiencia en etiquetas (\ref{sec:label-eff})
  & Datos anotados necesarios
  & PanDerm Large
  & HAM10000, PAD-UFES-20 \\
Segmentación (\ref{sec:segmentacion})
  & Localización de la lesión; arquitectura, dominio y tamaño
  & $5$ modelos SAM/SAM2 y controles
  & ISIC2018 (transferencia a ISIC2017, PH2) \\
Zero-shot multimodal (\ref{sec:zs})
  & Transferencia guiada por texto sin entrenamiento
  & DermLIP, BiomedCLIP, SigLIP
  & HAM10000, HIBA, MSKCC, DDI, Derm7pt clínico \\
Modelos generalistas (\ref{sec:llm})
  & ¿Compiten con los encoders especializados?
  & GPT-4o, Gemini~2.5~Pro, MedGemma 4B/27B
  & HAM10000 \\
Equidad por fototipo (\ref{sec:equidad})
  & Estabilidad del rendimiento entre fototipos
  & $5$ encoders
  & Fitzpatrick17k \\
Auditoría de solapamiento (\ref{sec:auditoria})
  & Fuga de información con Derm1M
  & Cómputo de \emph{hash}
  & $12$ externos \\
\hline
\end{tabular}
\end{table}

\subsection{Sondeo lineal (LP)}
\label{sec:lp}

El sondeo lineal mide el valor de las representaciones de un encoder
sin modificarlo: se congela el encoder y se entrena un clasificador
lineal sobre los \emph{embeddings} que extrae. La extracción aplica
la normalización canónica de cada modelo —parámetros ImageNet para
PanDerm y DINOv2, parámetros OpenAI CLIP para DermLIP, BiomedCLIP y
SigLIP— y los vectores se persisten en disco para no recomputarlos
entre experimentos. El clasificador es una regresión logística
multinomial entrenada sobre las particiones canónicas de la
tabla~\ref{tab:datasets-resumen}; su configuración exacta y los
parámetros de normalización por modelo figuran en el repositorio
reproducible (\ref{cap:repo}, sección~\ref{sec:repo-reproducibilidad}).
El protocolo se aplica a PanDerm Base y Large, DINOv2,
ConvNeXt-Large, EfficientNetV2-Large y la rama visual de DermLIP v2,
BiomedCLIP y SigLIP-Large.

\subsection{Ajuste fino supervisado (FT)}
\label{sec:ft}

A diferencia del sondeo lineal, el ajuste fino sí modifica el
encoder: entrena de forma conjunta el encoder visual y una cabeza de
clasificación. La receta reproduce la del paper original de
PanDerm~\cite{panderm2025}, para que la comparación sea directa:
optimizador AdamW con \emph{cosine annealing} y \emph{warmup}
lineal, regularización combinada (\emph{layer decay},
\emph{drop-path}, \emph{label smoothing}, Mixup y CutMix) y un
\emph{weighted random sampler} con pesos inversos a la frecuencia de
clase, sobre $50$ épocas en precisión mixta canónica con semillas
fijadas a $0$. El listado completo de hiperparámetros figura en el
repositorio reproducible (\ref{cap:repo},
sección~\ref{sec:repo-reproducibilidad}). En la evaluación final
sobre HAM10000 se añade \emph{test-time augmentation} (TTA): se
promedian los logits de $5$ augmentaciones deterministas por imagen.
El protocolo se aplica a PanDerm Base y Large sobre HAM10000,
PAD-UFES-20, DDI y MSKCC.

\subsection{Eficiencia en etiquetas}
\label{sec:label-eff}

Este protocolo responde a una pregunta práctica: cuántas etiquetas
hacen falta. Manteniendo el encoder congelado, se entrenan
clasificadores lineales (misma configuración que el sondeo lineal,
sección~\ref{sec:lp}) sobre submuestreos estratificados del
conjunto de entrenamiento de HAM10000 y PAD-UFES-20 al $1\%$,
$5\%$, $10\%$, $20\%$, $50\%$ y $100\%$. La
estratificación preserva la proporción de clases en cada subconjunto
y los submuestreos se generan con \texttt{random\_state}~$=42$.
PanDerm Large es el encoder principal; PanDerm Base se usa como
referencia adicional sobre HAM10000.

\subsection{Segmentación binaria de lesión}
\label{sec:segmentacion}

La segmentación evalúa la capacidad de localizar la lesión, y permite
separar la contribución de la arquitectura, del preentrenamiento médico
y del tamaño del \emph{backbone} en la familia \emph{Segment Anything}.
La evaluación se organiza en dos protocolos complementarios y varios
análisis de robustez.

\paragraph{Adaptación de SAM2.1-Large al dominio dermatológico.}
El primer protocolo adapta SAM2.1-Large mediante \emph{fine-tuning}
denso de su decodificador: se ajustan el decodificador de máscaras y
el \emph{promptable encoder}, mientras que el \emph{image encoder}
(Hiera-Large preentrenado) permanece congelado, de modo que solo
$\sim 4{,}2$\,M parámetros resultan entrenables ($\sim 17$\,MB de
\emph{checkpoint}). La optimización (AdamW, pérdida combinada
\emph{cross-entropy}$+$\emph{Dice}, normalización uniforme según la
convención de SAM2) se detalla en el repositorio reproducible
(\ref{cap:repo}, sección~\ref{sec:repo-reproducibilidad}). El
entrenamiento es sobre ISIC2018~\cite{isic2018} y la transferencia
fuera de dominio se evalúa sobre ISIC2017 y PH2 sin reentrenamiento.
Como referencias no adaptadas se incluyen SAM2 \emph{zero-shot} sin
\emph{fine-tuning} y un \emph{Convolutional AutoEncoder} (CAE-seg)
entrenado desde cero con idéntica pérdida y optimización.

\paragraph{Comparación controlada entre arquitecturas de
segmentación.}
El segundo protocolo separa los tres factores que determinan el
rendimiento ---generación de arquitectura, preentrenamiento de
dominio y tamaño de \emph{backbone}--- mediante una comparativa
pareada de cinco arquitecturas de segmentación \emph{promptable}:
MedSAM2-tiny~\cite{medsam2}, SAM2.1-tiny y SAM2.1-Large (familia
SAM2~\cite{sam2}), y SAM-Med2D~\cite{sammed2d} y MedSAM
v1~\cite{medsam} (familia SAM~\cite{sam2023}). Todas se entrenan
bajo un régimen común (\emph{fine-tuning} denso del decodificador y
el \emph{promptable encoder}, encoder visual congelado, \emph{prompt}
de \emph{bounding box} derivada de la máscara de referencia) sobre un
dataset dermatoscópico ampliado y deduplicado por hash MD5 de píxel
($14\,201$ imágenes únicas de ISIC2018, ISIC2017, HAM10000-seg y
PH2), y se evalúan sobre el mismo test canónico de ISIC2018
($N = 1\,000$, blindado frente al entrenamiento). La significación
de las diferencias de Dice se contrasta con el test de Wilcoxon
pareado~\cite{wilcoxon1945}.

\paragraph{Análisis complementarios.}
Para interpretar las diferencias observadas entre arquitecturas se
realizan cuatro análisis complementarios. Un control automático sin
caja ---una U-Net con codificador ResNet-101 preentrenado en
ImageNet--- se evalúa sin recibir \emph{bounding box} para aislar la
contribución del \emph{prompt}. La estabilidad de la mejor
configuración se estima por validación cruzada de cinco pliegues
estratificados por dataset de origen, junto a una ablación del rango
de adaptación de bajo rango (\emph{LoRA}, $r\in\{4,8,16,32\}$); su
robustez frente a la calidad del \emph{prompt} se mide bajo
perturbaciones controladas de la \emph{bounding box}, y la confianza
de las máscaras se calibra \emph{a posteriori} mediante
\emph{temperature scaling}. Por último, el techo de la métrica se
contextualiza frente al acuerdo inter-anotador del propio dataset.

\subsection{Clasificación zero-shot multimodal}
\label{sec:zs}

La clasificación zero-shot mide qué pueden clasificar los modelos
vision-lenguaje sin entrenamiento específico sobre las clases objetivo,
aprovechando su espacio compartido imagen-texto. Cada clase candidata se
formula como una o varias frases de plantilla, que el encoder textual
codifica y ---cuando hay varias por clase--- promedia en un único
vector; la imagen se codifica con el encoder visual y gana la clase de
mayor similitud coseno. Sobre HAM10000 se comparan tres formulaciones de
\emph{prompt}: (i)~genéricos en inglés, una plantilla por clase
(p.~ej.~\emph{``a dermoscopic image of a melanoma''});
(ii)~el conjunto A3 del paper de Derm1M~\cite{derm1m}, que sustituye el
nombre simple de cada enfermedad por varias descripciones clínicas y
visuales más informativas; y (iii)~una versión enriquecida de A3 con
descriptores clínicos (modalidad, localización, criterios visuales).
Para caracterizar la sensibilidad al tokenizador, el protocolo se
ejecuta con los dos tokenizadores de la familia DermLIP: PubMedBERT
a $256$ tokens (v1/v2) y GPT-2 con tokenizador de CLIP a $77$
tokens (original). Sobre HIBA, MSKCC, DDI y Derm7pt clínico se corre
la versión binaria melanoma/no-melanoma con \emph{prompts}
genéricos y A3. De este modo, el protocolo permite aislar el efecto
de tres componentes fundamentales del rendimiento zero-shot: el
modelo multimodal utilizado, la formulación textual de las clases y
el tokenizador empleado por la rama lingüística.

\subsection{Comparación con modelos de lenguaje multimodales}
\label{sec:llm}

Este protocolo evalúa si los grandes modelos multimodales de
propósito general, capaces de procesar simultáneamente imagen y
texto, compiten con los encoders especializados de las secciones
anteriores bajo un escenario común de clasificación diagnóstica. La
comparación contrapone dos paradigmas ---la especialización
dermatológica explícita frente a la escala de los modelos
generalistas---, con los modelos biomédicos generales (MedGemma) en
posición intermedia, lo que permite contrastar la contribución
relativa de la especialización de dominio frente a la escala del
preentrenamiento.

Para situar a los encoders frente a los grandes modelos
generalistas, el protocolo evalúa GPT-4o (API OpenAI),
Gemini~2.5~Pro (API Google), MedGemma~4B~Vision (local) y
MedGemma~27B Text (local en BF16) sobre el mismo \emph{split} de
test de HAM10000 ($1\,232$ imágenes) usado en sondeo lineal y
ajuste fino. El \emph{prompt} clínico estandarizado (unas $300$
palabras) plantea la tarea como apoyo al diagnóstico —no
sustitución del juicio clínico—, presenta las siete categorías de
HAM10000 con una descripción breve y pide el diagnóstico más
probable con razonamiento estructurado. La etiqueta se extrae por
\emph{parsing} determinista que mapea el nombre del diagnóstico al
código de HAM10000; las respuestas no mapeables (menos del $1\%$
en todos los modelos) cuentan como incorrectas. El modo principal es
\emph{zero-shot}; para MedGemma~27B se evalúa además una variante
ajustada con LoRA ($r = 8$) sobre el entrenamiento de HAM10000.

\subsection{Evaluación de equidad por fototipo}
\label{sec:equidad}

Diversos estudios han señalado que muchos conjuntos de entrenamiento
presentan una representación desigual de los fototipos cutáneos, lo
que puede traducirse en diferencias de rendimiento entre grupos
poblacionales~\cite{ddi}. Por ello, además de la precisión global,
este protocolo mide si el rendimiento se mantiene estable a lo largo
del espectro de fototipos de la escala de Fitzpatrick. Cinco
encoders (PanDerm Large, PanDerm Base, la rama
visual de DermLIP v2, DINOv2 ViT-L/14 y BiomedCLIP) se evalúan sobre
Fitzpatrick17k completo ($16\,577$ imágenes, $114$ patologías,
$6$ fototipos I-VI) por sondeo lineal (configuración de
sección~\ref{sec:lp}, con la normalización canónica de cada
modelo). Se reportan dos formulaciones: clasificación a $114$
patologías (grano fino) y a $3$ clases de malignidad
(benigno / no neoplásico / maligno), esta última habitual en
estudios de equidad sobre Fitzpatrick17k. Las métricas se calculan
por subgrupo sobre los seis fototipos, y la equidad se resume con la
brecha entre los extremos
$\mathrm{Gap}_{\mathrm{I}\to\mathrm{VI}} =
\mathrm{BAcc}_{\mathrm{I}} - \mathrm{BAcc}_{\mathrm{VI}}$. La
brecha entre fototipos extremos resume de forma sencilla la
estabilidad del rendimiento entre grupos; su interpretación, no
obstante, debe realizarse conjuntamente con la exactitud global, ya
que una diferencia reducida puede reflejar tanto un comportamiento
equitativo como un rendimiento uniformemente bajo en todos los
subgrupos. El experimento se verifica re-ejecutando el sondeo sobre
PanDerm Large y comprobando coincidencia hasta cuatro decimales con
el resultado original.

\subsection{Auditoría de solapamiento con Derm1M}
\label{sec:auditoria}

La interpretación de los resultados zero-shot requiere verificar
previamente la independencia entre los conjuntos de evaluación y los
datos utilizados durante el preentrenamiento. Si un modelo ha visto
durante su entrenamiento imágenes posteriormente utilizadas como
test, las métricas obtenidas pueden reflejar memorización parcial
del conjunto de evaluación y no capacidad real de generalización.
Por este motivo, se realiza una auditoría sistemática de
solapamiento entre los datasets evaluados y Derm1M, principal
dataset de preentrenamiento de la familia PanDerm y DermLIP.

La auditoría lo cuantifica: cada imagen se reduce
a su hash MD5 (biblioteca estándar de Python) y la intersección de
hashes identifica las imágenes presentes en ambos datasets. La
cardinalidad se reporta por dataset en la
sección~\ref{sec:auditoria-resultados}. Se auditan los doce datasets
externos, salvo SkinCon (no evaluado sobre clases diagnósticas).

El alcance de esta auditoría se declara con precisión. El hash MD5
detecta coincidencias \emph{exactas} a nivel de bytes y, por tanto,
identifica duplicados idénticos, pero no \emph{near-duplicates}
---recortes, reescalados, recompresiones o alteraciones leves de una
misma imagen---. La auditoría sostiene, en consecuencia, la ausencia
de \emph{solapamiento exacto}, no la ausencia completa de
contaminación visual. Esta elección prioriza la trazabilidad y el bajo
coste del hash exacto frente a la cobertura perceptual; el
\emph{perceptual hashing} o la comparación por \emph{embeddings}
visuales constituyen su extensión natural, y las implicaciones de este
alcance acotado para la interpretación de los resultados se retoman en
la sección~\ref{sec:lim-leakage}.
